\section{Related Work}
\label{sec:related_work}

We organize learning-based colorization into four paradigms that emerged in
succession, each addressing a limitation of the previous one. This section
gives the historical and conceptual map; the architectural details of every
model we evaluate appear in Sec.~\ref{sec:method}.
Table~\ref{tab:taxonomy} summarizes the taxonomy and maps each paradigm to the
concrete model we benchmark.

\subsection{CNN classification (2016)}
Zhang \emph{et al.}~\cite{zhang2016colorful} treat colorization as per-pixel
classification over a quantized $ab$ palette rather than regression, explicitly
modeling color multimodality and rebalancing the loss toward rare, saturated
colors. The result was a step change in perceptual vividness and the first
strong fully-automatic colorizer. This is the model we reimplement
(Sec.~\ref{sec:method:zhang16}); its design is also the conceptual backbone
that the next three paradigms either extend or react against.

\subsection{Interactive CNN (2017)}
The follow-up of Zhang \emph{et al.}~\cite{zhang2017realtime} adds a second
input branch for sparse user hints (clicked color points) and U-shape skip
connections, fusing learned priors with user intent in real time. The same
network can run with \emph{zero} hints, in which case it behaves as a purely
automatic colorizer---the mode in which we benchmark it
(Sec.~\ref{sec:method:zhang17}) for a fair comparison.

\subsection{GAN (2020)}
Adversarial training replaces a hand-designed loss with a learned
discriminator that rewards realism. ChromaGAN~\cite{vitoria2020chromagan}
couples an adversarial objective with a semantic class-distribution prior; its
contemporaneous, most widely-deployed practical descendant is
DeOldify~\cite{antic2019deoldify}, whose \emph{NoGAN} recipe stabilizes
training and produces robust, restoration-grade colorizations
(Sec.~\ref{sec:method:deoldify}). We use it as a SOTA-class anchor.

\subsection{Diffusion (2022)}
Conditional diffusion models cast colorization as iterative denoising.
Palette~\cite{saharia2022palette} frames it as a general image-to-image
diffusion task, while ControlNet~\cite{zhang2023controlnet} adds spatial
conditioning to a pretrained text-to-image backbone via a trainable encoder
copy with zero-convolution adapters (Sec.~\ref{sec:method:controlnet}). We
attempt a ControlNet pipeline as the diffusion representative;
Sec.~\ref{sec:results} reports the reproducibility obstacle we encountered.

Evaluation in all four paradigms increasingly relies on learned perceptual
metrics; we adopt LPIPS~\cite{zhang2018lpips} alongside PSNR and SSIM, and
benchmark on COCO~2017~\cite{lin2014coco}.

\begin{table}[t]
\centering
\caption{Four colorization paradigms and the models evaluated in this study.}
\label{tab:taxonomy}
\begin{tabular}{@{}llll@{}}
\toprule
Paradigm & Representative & Year & Our proxy model \\
\midrule
CNN              & Zhang~\cite{zhang2016colorful}      & 2016 & Zhang16 (ours) \\
Interactive CNN  & Zhang~\cite{zhang2017realtime}      & 2017 & Zhang17 (auto) \\
GAN              & ChromaGAN~\cite{vitoria2020chromagan} & 2020 & DeOldify \\
Diffusion        & Palette~\cite{saharia2022palette}   & 2022 & ControlNet \\
\bottomrule
\end{tabular}
\end{table}
